\documentclass[10pt]{article}
\usepackage[margin=20mm]{geometry}
\usepackage{amsmath,amssymb,mathtools,fontspec,unicode-math,xurl,hyperref}
\IfFontExistsTF{DejaVu Serif}{\setmainfont{DejaVu Serif}}{\setmainfont{Latin Modern Roman}}
\setmathfont{latinmodern-math.otf}
\setlength{\parindent}{0pt}\setlength{\parskip}{6pt}\sloppy
\title{Initial verification of
Locally Homogeneous Axiom A Flows I:
Projective Anosov Subgroups and Exponential Mixing}
\author{}
\date{Math 2025 initial review — M25-0384}
\begin{document}
\maketitle
Abstract: One supported Major issue concerns the claimed entropy-valued leading resonance of the scalar transport operator. A compact surface-lattice example satisfies the stated hypotheses, but its scalar resolvent is holomorphic throughout the open right half-plane. Four supported Minor issues concern the elementary-group boundary, a relative sign in the pseudo-hyperbolic embedding, resolvent conventions, and a local-product quantifier. A separate uniformity obligation in the non-integrability argument remains NeedsEvidence; no counterexample to the intended non-elementary exponential-mixing theorem is inferred from that obligation. The geometric construction, linear distortion argument, reference hypotheses, and application-specific time scales are checked with the qualifications below. The scientific decision remains pending. This initial review is not independent acceptance.\par
\section*{Identity, version, and review scope}
The assigned work is Benjamin Delarue, Daniel Monclair, and Andrew Sanders, Locally Homogeneous Axiom A Flows I: Projective Anosov Subgroups and Exponential Mixing, assigned DOI `10.1007/s00039-025-00712-2`. The supplied and reviewed source is `arXiv:2403.14257v1`, dated 21 March 2024, with 54 physical pages [paper]. All page references below are one-based physical pages; their printed numbers agree. Conclusions concern this supplied preprint, not an unexamined later version of record.\par
The source has 662,551 bytes and SHA256\par
b6f51ce40e01815eb090cb77fa5087f1
3aac0c66d6eca0a481e08e2394886112\par
The complete source text and all 54 source-page images were actually read or viewed in this scientific conversation. The title-page authors are used; the incorrect running-head surname does not change the paper’s identity. Exact reference versions, selected text pages, image paths, byte counts, and hashes are recorded in the evidence inventory. Acquiring or parsing a reference does not constitute reading it, and reading selected passages is not represented as reading an entire reference. The dispatch, task, complete prompt, source manifest, and frozen contract 0.7.1 remain bound in the provenance record.\par
\section*{Claims, hypotheses, and dependencies}
For a real vector space \(V\) of dimension \(d>1\), the paper uses \[L=\{[v:\alpha]\in\mathbb P(V\oplus V^*):\alpha(v)>0\},\qquad
 \phi^t[v:\alpha]=[e^t v:e^{-t}\alpha].\] The normalization \(\alpha(v)=1\) is a double cover of this projective space. On it, the Liouville form \(\tau=\alpha\,dv\) has Reeb field \((v,-\alpha)\). The primal and dual transverse distributions give the stable and unstable foliations.\par
Definition 2.14 uses a finitely generated hyperbolic group and a uniform eigenvalue-gap bound \[\lambda_1(\gamma)-\lambda_2(\gamma)
 \geq c|\gamma|_\infty-c',\qquad c,c'>0,\] where \(\lambda_j\) are logarithms of eigenvalue moduli and \(|\gamma|_\infty\) is stable word length. Torsion-freeness is separately needed for the free quotient. This definition does not itself exclude an infinite cyclic group.\par
Theorems A and B construct a flow-invariant discontinuity domain, identify a compact invariant set with the full ordered-pair refraction flow, and assert a unique transitive basic set with analytic ambient stable and unstable distributions. The proof chain is: transverse limit maps and convergence dynamics; properness of the action; the Hopf parametrization and cocompactness; identification of non-wandering, periodic-closure, and trapped sets; then the Anosov bundle estimates. Compactness of \(\mathcal K_\Gamma\) does not mean that \(\mathcal M_\Gamma\) is compact.\par
For mixing, Section 4 explicitly adds non-elementarity, and the principal theorem adds irreducibility. The linear pseudo-exponential map, spanning of the infinitesimal limit set, and a finite compact cover give regular distortion. The other two inputs to Stoyanov’s theorem are Lipschitz stable holonomy and local non-integrability. Theorem C/Theorem 4 concerns every Hölder potential and Hölder observable, with constants permitted to depend on the potential and Hölder exponent. The analytic flow supplies the differentiability required in Stoyanov’s stated mixing corollary [stoyanov].\par
Section 5 returns to the broader non-trivial torsion-free scope. Theorems D–E and 5–8 concern Euler products over primitive conjugacy classes, meromorphic continuation, a leading pole, a zero-free strip, and prime-orbit asymptotics. Theorems F and 9–11 separately discuss the scalar operator \(A=-X+U\), localized resolvents, resonant states, and wavefront-compatible products. The pseudo-hyperbolic and Benoist–Hilbert applications in Theorems H/G and 12/13 must retain the required group range, operator type, and time normalization. These distinctions are material to the findings below.\par
\section*{Problem 1: positive entropy is not a scalar leading resonance}
Severity: Major. Location: Theorem F, p. 7, especially the \(U=0\) assertions; the scalar definition and resolvent on p. 41; Theorem 12(3), p. 47, and Theorem 13(3), p. 51. Section 5.2.2 on p. 43 must be reconciled with these statements.\par
\subsection*{An admissible compact example}
Choose a closed oriented hyperbolic surface of genus at least two. Lift its torsion-free cocompact holonomy in \(\mathrm{PSL}_2(\mathbb R)\) to a faithful subgroup \(\Gamma<\mathrm{SL}_2(\mathbb R)\) not containing \(-I\). Such lifts exist: the obstruction is the Euler class modulo two, and the surface Euler number \(2g-2\) is even. Projection to the surface group is injective, so the lifted group is torsion-free and cocompact modulo the finite center. Its standard action on \(V=\mathbb R^2\) is irreducible, since a cocompact non-elementary Fuchsian group has no invariant real line. It is projective Anosov: the eigenvalue gap is the hyperbolic translation length, uniformly comparable to stable word length by the orbit quasi-isometry. Its limit set is all of \(\mathbb P(\mathbb R^2)\). Thus the example satisfies the theorem’s group, dimension, compactness, smoothness, and potential hypotheses. The paper itself explicitly includes uniform \(\mathrm{SL}_2(\mathbb R)\) lattices among its compact cases on p. 3.\par
Here \(L\) is naturally \(\mathrm{PSL}_2(\mathbb R)\). For a normalized pair \(\alpha(v)=1\), there is a unique \(w\in\ker\alpha\) with \(\det(v,w)=1\). The matrix with columns \(v,w\) identifies the normalized double cover with \(\mathrm{SL}_2(\mathbb R)\); \((v,\alpha)\sim(-v,-\alpha)\) gives the quotient by \(\{\pm I\}\). The flow becomes right multiplication by \(\operatorname{diag}(e^t,e^{-t})\). Because the limit set is full, every transverse pair belongs to the domain and to the lifted invariant set. Consequently \[\mathcal M_\Gamma=\mathcal K_\Gamma
 =\overline\Gamma\backslash\mathrm{PSL}_2(\mathbb R)\] is compact. This is the geodesic flow with twice the usual curvature-\(-1\) time parameter, so its entropy \(h\) is positive (\(h=2\) with this normalization).\par
\subsection*{Direct resolvent exclusion}
Equation reference: eq:scalar-resolvent. Take \(U=0\), which satisfies every support requirement, including on a compact ambient manifold. The contact volume \(\mu=\tau\wedge d\tau\) is finite and invariant. Hence \[T_t f=f\circ\phi^{-t}\] is a unitary group on \(L^2(\mu)\), with generator \(A=-X\). For every \(\operatorname{Re}\lambda>0\), \[
 (A-\lambda)^{-1}
 =-\int_0^\infty e^{-\lambda t}T_t\,dt,\qquad
 \|(A-\lambda)^{-1}\|\leq\frac{1}{\operatorname{Re}\lambda}.\] The strong Laplace integral has tails converging in operator norm, locally uniformly on the open right half-plane. Differentiation is controlled by the integrable bounds \(t^k e^{-t\operatorname{Re}\lambda}\), so it defines a holomorphic bounded-operator family there. Composing with the continuous inclusions from smooth functions to \(L^2\) and from \(L^2\) to distributions gives the same conclusion for the distribution-valued resolvent. Its unique meromorphic continuation cannot have a pole at the positive real number \(h\). On the constant function, \[(A-\lambda)^{-1}1=-\lambda^{-1}1.\] Thus zero is the scalar resonance exhibited by constants, whereas the asserted positive-entropy scalar resonance is absent. This is an analytic contradiction, not a numerical failure to detect a pole.\par
The counterexample also explains the error’s nature. Humbert’s original Theorem 1 distinguishes the scalar critical pressure \(P(U+J^u)\) from the pressure \(P(U)\) for the appropriate degree of differential forms, where \(J^u\) is the negative infinitesimal unstable Jacobian [humbert]. In the compact constant-curvature example, \(J^u=-2\), \(P(0)=2\), and \(P(J^u)=0\). Dyatlov–Guillarmou likewise distinguish the scalar zero resonance from the entropy-related form contribution to the zeta function [dg2016].\par
\subsection*{Impact and repair boundary}
The failure changes a headline spectral assertion and the claimed entropy-state identification for the specified scalar operator. Changing only the overall sign of a resolvent does not move its poles and therefore cannot repair it.\par
The paper’s p. 43 does mention lifts to differential forms. That correct observation is acknowledged here; it does not make the expressly scalar Theorem F true. A repair must specify the operator, form degree, orientation twist when needed, pairing, and resulting invariant distribution, then prove the claimed open-basic-set version. Alternatively a change of potential changes the operator and requires its own regularity and pressure analysis. The compact Anosov result alone does not certify every open-set extension.\par
The same positive-entropy scalar claim fails for a closed hyperbolic geodesic flow in the applications. Moreover, a Hölder conjugacy alone does not define pullback of arbitrary distributions with specified wavefront sets; the low-regularity Hilbert-flow statement needs an explicit functional interpretation. These are repair obligations, not a claim that exponential mixing, zeta continuation, or the geometric construction is refuted by this compact example.\par
\section*{Problem 2: the cyclic boundary has two periodic components}
Severity: Minor; range correction. Location: Theorems A/B, pp. 3–4; Theorem D, p. 6; Proposition 2.9 and Definition 2.14, pp. 13–14; Theorem 2 and Lemma 3.12, pp. 22–23; the elementary-case sentence on p. 25; Theorem 6, p. 39. The same omitted boundary restriction affects the literal one-dimensional applications.\par
Let \[g=\begin{pmatrix}2&0\\0&1/2\end{pmatrix},\qquad
 \Gamma=\langle g\rangle\simeq\mathbb Z,\qquad \ell=\log2.\] This is a non-trivial torsion-free hyperbolic group. With the standard generator, \(|g^n|_\infty=|n|\), and \[\lambda_1(g^n)-\lambda_2(g^n)=2|n|\log2.\] Thus the literal projective Anosov gap condition holds with positive constants; for example \(c=c'=\log2\) works for every integer \(n\). The primal boundary points are \(\xi(+)=\mathbb Re_1,\ \xi(-)=\mathbb Re_2\), and the dual ones are \(\xi^*(+)=\mathbb Re_2^*,\ \xi^*(-)=\mathbb Re_1^*\). They satisfy the required transverse-pair condition.\par
The ordered boundary-pair space has two points, \((+,-)\) and \((-,+)\), both fixed by \(\Gamma\). Each corresponding real fiber quotients to a circle. The translations are respectively \(+\ell\) and \(-\ell\), so the canonical \(\mathcal K_\Gamma\) and the full refraction flow are two disjoint periodic circles, each of period \(\ell\). No orbit is dense in their union. This contradicts the canonical construction’s transitivity and unique-basic-set assertion, and the joint A/B identification with the full refraction flow. It does not prove that no different domain could be chosen to retain a single periodic component; Theorem A’s existential wording should not be overstated.\par
Equation reference: eq:cyclic-zeta. The only primitive conjugacy classes in \(\mathbb Z\) are \(g\) and \(g^{-1}\). They are distinct, and both have period \(\ell\). Consequently the exact Euler product is \[
 \zeta_\Gamma(s)=\bigl(1-e^{-s\ell}\bigr)^{-2}.\] The flow entropy is zero, and \(s^2\zeta_\Gamma(s)\to\ell^{-2}\). There is a double pole at zero, contradicting the simple-pole assertion for this allowed case. The independent rational series check gives \[\frac{z^2}{(1-e^{-z})^2}
 =1+z+\frac{5}{12}z^2+\frac1{12}z^3+\frac1{240}z^4+O(z^5),
 \qquad z=s\ell .\]\par
The interval \(C=\{[x:y]:x,y>0\}\subset\mathbb P(\mathbb R^2)\) is properly strictly convex under Definition 6.10. The action multiplies \(x/y\) by \(4\), and \(C/\Gamma\) is a compact Hilbert circle of length \(\ell\). Its unit sphere bundle again has two oriented periodic circles. Periodic rotation is not mixing; a nonconstant Fourier mode on either circle has correlations that do not decay. Thus Theorem G/13 needs a non-elementary or dimension restriction under the literal stated definition. The analogous \(p=1,q=0\) pseudo-hyperbolic example is also permitted by the stated range; Theorem H omits the irreducibility condition retained in Theorem 12(1).\par
The repair is to state non-elementarity wherever transitivity, positive entropy, mixing, or a simple leading pole is used, or to give the elementary case separately with its two-orbit product. BCLS explicitly excludes virtually cyclic groups in its standing conventions on p. 9 [bcls]. Section 4 of the reviewed paper already excludes elementary groups, and the diagonal example is reducible, so it does not refute the intended irreducible mixing theorem. The finding is a localized but genuine scope error.\par
\section*{Problem 3: the pseudo-hyperbolic embedding has the wrong sign}
Severity: Minor. Location: the positive-pairing target on p. 44 and the displayed map \(\Phi_\partial^{p,q+1}\) on p. 45.\par
Write \(B=\langle\,\cdot,\cdot\,\rangle_{p,q+1}\). A unit spacelike tangent lift satisfies \[B(x,x)=-1,\qquad B(v,v)=1,\qquad B(x,v)=0.\] The printed map is \([x:v]\mapsto[x+v:x-v]\), while its target requires \(B(v_1,v_2)>0\). Directly, \[B(x+v,x-v)=B(x,x)-B(v,v)=-2.\] The sign is unchanged by common projective rescaling. Hence the displayed formula does not land in its stated target for any normalized input. For a completely concrete permitted case take \(p=2,q=0\), \(B=\operatorname{diag}(1,1,-1)\), \(x=(0,0,1)\), and \(v=(1,0,0)\). Every domain condition holds.\par
One consistent repair is \[[x:v]\longmapsto[x+v:v-x],\qquad
 [v_1:v_2]\longmapsto[v_1-v_2:v_1+v_2].\] The second formula is homogeneous. For null \(v_1,v_2\) with \(B(v_1,v_2)>0\), a common normalization makes the resulting pair unit tangent. The geodesic equations \(x_t=\cosh(t)x+\sinh(t)v,\ 
v_t=\sinh(t)x+\cosh(t)v\) give \(x_t+v_t=e^t(x+v)\) and \(v_t-x_t=e^{-t}(v-x)\), exactly the required flow relation. The independent calculation verifies these identities at the nontrivial rational point \(\cosh t=5/4,\ \sinh t=3/4\), together with the general bilinear identities used above.\par
This sign error invalidates the map as printed. The corrected relative sign supplies the intended equivariant local construction. All subsequent domain and inverse formulas should be made consistent with that correction and their common normalization. No counterexample to the intended geometric mixing conclusion is claimed from this repairable formula.\par
\section*{Problem 4: resolvent residue and starting half-plane conventions}
Severity: Minor. Location: Section 5.2 and Theorem 9, p. 41. The quoted spectral-parameter convention in Theorem 10, p. 42, also requires attention.\par
Keep the compact admissible example from Problem 1 and choose the smooth real potential \(U=c>0\). Its support is compact because the ambient manifold is compact. On constants, \[A1=c1,\qquad (A-\lambda)^{-1}1=(c-\lambda)^{-1}1.\] The residue at \(c\) sends \(1\) to \(-1\). Thus, for the printed definition \(\Pi_c=\operatorname{Res}_{\lambda=c}(A-\lambda)^{-1}\), \[(\Pi_c^2-\Pi_c)1=2\,1\ne0.\] It is not a projection. The required convention is \(-\operatorname{Res}(A-\lambda)^{-1}\), or instead the resolvent \((\lambda-A)^{-1}\) with its positive residue. The rank and pole locations survive this sign correction, but the projector and normalization statements must use one convention consistently.\par
The same example refutes the assertion that the initial \(L^2\)-resolvent exists for every \(\operatorname{Re}\lambda>0\) for arbitrary \(U\): it has a pole at \(c>0\). A sufficient initial half-plane for the transport integral with invariant volume is \(\operatorname{Re}\lambda>\sup\operatorname{Re}U\). The potential-dependent starting region and the resolvent sign are explicit in the original open-system framework [dg2016].\par
Jin–Tao use \(P=-iX\) and strips \(\operatorname{Im}z>-\beta\) [jintao]. For the present \(A=-X=-iP\), the corresponding variable is \(\lambda=-iz\), so \(\operatorname{Re}\lambda=\operatorname{Im}z\). A direct transfer of their strip must therefore be written \(\operatorname{Re}\lambda>-\beta\). Their open-system statements, rather than an unqualified use of the compact global statement, are the relevant input after localization. This identifies a citation and convention mismatch; it is not offered as a separate counterexample to every counting inequality printed in Theorem 10. None of these local convention repairs resolves Problem 1.\par
\section*{Problem 5: the local-product claim needs a basic-set restriction}
Severity: Minor. Location: Lemma 4.3(3), p. 28; compare Definition 4.4, p. 29, where subsequent uses already take base points in \(\mathcal K_\Gamma\).\par
The lemma states that for every ambient point \(y\in B(x,2\varepsilon_1)\), the local intersection \(W^s_{\varepsilon_0}(x)\cap W^{cu}_{\varepsilon_0}(y)\) belongs to \(\mathcal K_\Gamma\). Consider a torsion-free classical Schottky group in \(\mathrm{SL}_2(\mathbb R)\). It is non-elementary, irreducible in the standard representation, convex cocompact, and projective Anosov. Its projective limit set is a Cantor set. In the paper’s model, \[\mathcal K_\Gamma\cap W^s_{\mathrm{loc}}(x)\] has a Cantor transverse parameter: the stable leaf fixes the primal endpoint and varies the dual endpoint, which must lie in the dual limit set for membership in \(\mathcal K_\Gamma\).\par
For every sufficiently small neighborhood of \(x\), choose \(y\in W^s_{\varepsilon_0}(x)\setminus\mathcal K_\Gamma\) in that neighborhood. Such points exist because the dual Cantor set contains no interval. Also \(y\in W^{cu}_{\varepsilon_0}(y)\). The local stable/central-unstable intersection, which is unique by the ordinary ambient local-product property, is therefore \(y\notin\mathcal K_\Gamma\). This satisfies the lemma’s stated ambient quantifier and contradicts its conclusion.\par
The local-product property inside a locally maximal hyperbolic set requires \(y\in\mathcal K_\Gamma\) as well. Replacing the ambient quantifier by \(y\in\mathcal K_\Gamma\cap B(x,2\varepsilon_1)\) repairs this step. The subsequent holonomy arguments use points in the basic set, so this finding does not refute their intended conclusions. Likewise, the parallel transport mentioned in item (5) should be restricted to the leaf along which it was defined.\par
\section*{NeedsEvidence: a uniform cone angle in non-integrability}
Identifier: NE001. Location: Definition 4.19, p. 34; Lemma 4.20 and Lemma 4.21, pp. 35–36; application in Theorem 3, p. 37.\par
The strong local non-integrability condition has the quantifier order \[\exists d_0>0\quad
 \forall\varepsilon\in(0,\varepsilon_0)\quad
 \forall z\in\mathcal K_\Gamma\quad
 \forall w\in E^u_z,\ \|w\|=1\quad\cdots.\] Lemma 4.20 proves a cone-angle bound after fixing a neighborhood scale \(\delta>0\). Its compactness argument is uniform in the base point and unit vector for that fixed scale. In Lemma 4.21, however, \(\varepsilon\) is chosen first, then \(\delta\) small enough to fit the stable neighborhood of radius \(\varepsilon\), and then \(d_0\) is supplied by Lemma 4.20. The argument as written produces \(d_0=d_0(\delta(\varepsilon))\); it does not establish a positive lower bound independent of \(\varepsilon\).\par
The logarithmic estimate itself is valid at a fixed scale: the displayed temporal separation is \(-\log(1+\beta(u)/\beta(v'))\), and for a sufficiently small argument its absolute value is bounded below by one half of that argument’s absolute value. The uniform leaf-chart and parallel transport estimates then give the stated positive \(\kappa\). This verifies the local calculation without silently exchanging the order of the quantifiers.\par
Stoyanov’s original LNIC also fixes its cone parameter before the small-neighborhood quantifier, but tests tangent directions to the basic set and permits a shifted base point and two stable points [stoyanov, p. 7]. It is weaker than the paper’s proposed strong condition. The requested evidence is either a uniform-angle argument for Definition 4.19 or a direct verification of that original weaker hypothesis. Irreducibility and spanning at each fixed scale alone do not supply the missing uniform lower bound in the written proof. No admissible counterexample to the mixing theorem has been established here; this item remains an unresolved proof obligation.\par
\section*{Proof-level coverage and independent counterchecks}
\subsection*{Convergence dynamics, properness, and the trapped set}
The boundary convergence statements, transverse primal/dual maps, stable word-length estimates, and the properness proof in Sections 2–3 were followed with the elementary exception in Problem 2 kept explicit. The eigenvector-growth contradiction in Theorem 1 uses the positive pairing: if the primal vector grows and a projective normalization shrinks to zero, the normalized pairing would tend to zero, contradicting positivity. The dual alternative is symmetric.\par
In Lemma 3.5 the dual attracting exponent should be \(-\lambda_d(\gamma)=\lambda_1(\gamma^{-1})\), rather than \(\lambda_1(\gamma)\) in general. Positivity of that exponent suffices for the first contradiction. The subsequent primal calculation still gives the claimed period \(T=\lambda_1(\gamma)\). This is a local proof correction.\par
The written proof of Lemma 3.10 has incompatible inverse and primal/dual labels. Its conclusion admits a direct correction. Use normalized lifts with \[w_k=e^{t_k}\delta_k v_k,\qquad
 \beta_k=e^{-t_k}\delta_k\alpha_k,
 \qquad t_k\longrightarrow+\infty.\] The sequence \(\delta_k\) is unbounded. After the fixed multiplication supplied by Lemma 2.11, put \(\gamma_k=f\delta_k\), with distinct boundary limits \(\eta_+,\eta_-\). Then \(\gamma_kv_k=e^{-t_k}fw_k\to0\), so the growth lemma shows that \([v]\) is not transverse to \(\xi^*(\eta_-)\). Membership of \([v:\alpha]\) in the domain forces \([\alpha]\) transverse to \(\xi(\eta_-)\). Uniform dual convergence gives \(f[\beta]=\xi^*(\eta_+)\). Since \(f\beta(fw)=1\), the line \(f[w]\) is transverse to this dual hyperplane. Applying inverse convergence to \(f[w_k]\) and using \(\gamma_k^{-1}fw_k=e^{t_k}v_k\) gives \([v]=\xi(\eta_-)\). Thus \([v]\in\Lambda_\Gamma\) and \([\beta]\in\Lambda_\Gamma^*\), as required. The original uniform convergence result and its proof are available in the checked DGK passages [dgk]. No failure of the trapped-set conclusion is inferred from the printed labels.\par
The Hopf coordinate and its logarithmic cocycle were checked directly. Cocompactness and positive time reparametrization are supported by the original Sambarino construction in its stated geometric setting and by the BCLS general-group construction [sambarino, bcls]. The latter’s non-elementary convention must be retained. Hyperbolicity in Lemma 3.13 uses contraction of the pulled-back projective tangent bundles, as in the actual Guichard–Wienhard definition [gw]. The ambient formula \(d\phi^t(w,0)=(e^tw,0)\) alone would not suffice; the paper correctly explains that distinction on pp. 24–25.\par
\subsection*{Regular distortion, holonomy, and mixing}
Equations (4.3)–(4.9) give the analytic leaf charts and rational holonomy formulas on their transverse domain. Restricting these smooth maps to a compact neighborhood yields the uniform Lipschitz bounds used later. The corrected basic-set quantifier of Problem 5 is enough for the stated applications.\par
The local spanning argument in Lemma 4.8 can be made precise by first choosing a smaller open subset whose closure avoids a repelling endpoint and then taking a sufficiently large contracting power. Finite-dimensional invariance and minimality of the non-elementary boundary action yield the global span. For Lemma 4.15 choose its basis points in the ball of radius \(\delta/2\), then shrink the neighborhood so their transported images remain in the ball of radius \(\delta\). If \(L_yu_i=e_i\), \(\|L_y\|\leq b\), and \(A\) is the relevant backward linear map, then for \(\|u\|\leq\varepsilon\), \[\|Au\|
 \leq\sqrt m\,b\,\varepsilon\max_i\|Au_i\|
 \leq\sqrt m\,b\,\varepsilon\,
       \operatorname{diam}A(\Lambda^u\cap B_\delta).\] A factor of two passes from radius to diameter. Compactness provides finitely many such charts. This verifies the regular-distortion mechanism without assuming a pinching condition absent from the paper.\par
The stronger distortion inequality implies both of Stoyanov’s distortion requirements: the comparison for fixed radii and the arbitrarily small diameter ratio obtained by shrinking the smaller radius. Stoyanov’s Corollary 1.5 explicitly allows Hölder potentials when the flow is \(C^5\); the present analytic flow meets that regularity. There is therefore no supported objection based simply on Hölder, rather than Lipschitz, potentials. The outstanding input is precisely the quantified non-integrability obligation NE001, not those already verified hypotheses.\par
\subsection*{Zeta functions and microlocal localization}
The primitive-class/period correspondence is compatible with \(\lambda_1(\gamma)\), including both orientations in the cyclic calculation. The global continuation input is genuinely local near a compact locally maximal hyperbolic set: the Dyatlov–Guillarmou afterword constructs a suitable strictly convex isolating block by modification away from the trapped set [afterword]. The noncompact ambient manifold is therefore not by itself an obstruction. The original open-system resolvent and localization-independence argument were checked [dg2016].\par
Borns-Weil–Shen use the orientation line bundle to handle nonorientable invariant bundles [bws]. Their alternating form trace identity accounts for the sign of the Poincaré determinant. Orientability should not be silently added as an assumption to the reviewed theorem. The wavefront restriction for products of resonant and coresonant states is present in the original theorem. The actual source images on pp. 7 and 42 contain the complex conjugations; losses of overbars in extracted text were not treated as mathematical errors.\par
For the unweighted counting conclusion, Stoyanov’s stated prime-orbit corollary supplies an absolute exponential remainder. Dividing by \(\operatorname{Li}(e^{ht})\) introduces a factor of order \(t\), which is absorbed by choosing a slightly weaker exponent. This is compatible with the existential constant in Theorem 8. The signs in the displayed remainders of Theorems 12(4) and 13(4) should agree with the decaying sign in Theorem 8; as printed their bounds are weak rather than the claimed decay.\par
The weighted continuation in Theorem 5 has been checked against Pollicott’s Theorem 2 and its proof, including the general Axiom A statement on printed p. 422 (physical p. 11 of the acquired scan) [pollicott]. The result concerns a weakly mixing flow on a non-trivial basic set and a Hölder potential. Exponential mixing supplies the required weak mixing and holomorphy of the correlation transforms. The proof uses a finite-to-one symbolic coding with a nonvanishing analytic correction factor, and spectral projections that detect the non-leading poles through suitable pairs of observables. Centering the correlations removes the leading pressure contribution. Holomorphy for one arbitrarily chosen pair alone would not suffice. The construction is local to the compact basic set; compactness of the ambient quotient is not required by the coding argument. The group-range issue and NE001 remain separate obligations.\par
The Dolgopyat–Pollicott author version was also read in full [dp1998]. Its pp. 4–6 describe a collapsed base in \(\mathbb R^n\), where \(n\) is the unstable dimension, and specialize to intervals for surfaces. The reviewed p. 40 imports this interval model into the general-dimensional setting and uses the exact point-evaluation formula (5.5). The survey itself prints such an equality without a detailed derivation; that line does not independently verify the claimed exact identity for the present general model. This proof sketch should be replaced by a justified general-dimensional periodic-sum estimate, or more directly by Stoyanov’s Corollary 1.4. The latter, actually checked above, supplies the desired unweighted zero-free continuation and counting conclusion once its Section 4 hypotheses are met. Thus a verified alternative dependency is available; no independent counterexample to Theorem 7 is asserted.\par
The germ terminology on p. 39 should also be distinguished from restriction to \(\mathcal K_\Gamma\); a smooth function can vanish on that set without vanishing in a neighborhood. The orbit weights only require the restriction.\par
\subsection*{Pseudo-hyperbolic and Hilbert applications}
The convex-cocompactness, negative-limit-set, and projective-Anosov hypotheses were compared with the actual DGK definitions and statements [dgk]. The repaired embedding in Problem 3 preserves the flow and the intended domain relation. Restriction to the trapped set is needed when the ambient spacelike geodesic flow is incomplete. The irreducibility qualifier in Theorem 12(1) should also appear in the corresponding introductory claim.\par
The adjoint representation has highest eigenvalue logarithm \(\lambda_1-\lambda_d\). The projective Anosov bounds for \(\gamma\) and \(\gamma^{-1}\) give the required two extreme gaps, consistently with the original representation-theoretic criterion [gw]. The symmetric boundary graph is Lipschitz by the checked convex-gradient argument in Glorieux–Monclair–Tholozan [gmt]: the monotonicity inequality makes its graph spacelike, and projection onto a maximal positive subspace gives a Lipschitz graph. A smooth-boundary claim is not needed.\par
In Lemma 6.17 the positive cones and trace-one representatives give the intended inverse to the rank-one tensor formula. The explicit calculation gives \[\Psi\circ\phi_{\rm BH}^t=\phi^{2t}\circ\Psi.\] Accordingly, \(T_{\rm Ad}=2T_{\rm BH}\), \(h_{\rm BH}=2h_{\rm Ad}\), and \(\zeta_{\rm BH}(s)=\zeta_{\rm Ad}(s/2)\). The Hilbert length formula in the original reference confirms the factor \(1/2\) [gmt]. Hölder conjugacy and a constant time change preserve the appropriate qualitative mixing conclusion with adjusted constants and exponents. They do not preserve the numerical period without rescaling, nor do they by themselves transport arbitrary microlocal distributional structures. The ellipsoid case is treated by the classical compact hyperbolic flow; the non-ellipsoid argument requires the stated Zariski-density/irreducibility input. The one-dimensional boundary remains Problem 2.\par
\section*{Evidence disposition and limits}
The supported disposition is one Major problem and four Minor problems. Each has an exact location, an admissible counterexample or direct analytic contradiction, an impact statement, and a stated repair boundary. NE001 is retained separately as a proof obligation. The written independent arguments establish the decisive contradictions; the native rational computation corroborates the displayed algebra and Laurent coefficient and does not substitute for those arguments.\par
The source was read in full. Thirteen original public reference files were checked at their recorded statement and proof pages. The Pollicott scan and the Dolgopyat–Pollicott author version were read in full; the other references were read at the recorded selected pages, not recursively through every foundational citation. The distinction between acquired preprint versions and later publication versions is retained. The repaired dependency for the unweighted zeta proof and the non-integrability uniformity issue remain visible in the coverage and findings records. No author code was run, and absence of author code is not classified as a mathematical defect.\par
The real compilation command, exit code, compiler and dependency bytes before and after, final PDF bytes, and actual page-inspection hashes are recorded outside this report so that its own PDF hash is not circularly embedded. A delivery flag is earned only after these records match the current files. The scientific decision is `pending`, with no human reviews entered. Initial delivery does not constitute independent acceptance.\par
99 B. Delarue, D. Monclair, and A. Sanders, Locally Homogeneous Axiom A Flows I: Projective Anosov Subgroups and Exponential Mixing, arXiv:2403.14257v1, 21 March 2024. Assigned DOI: 10.1007/s00039-025-00712-2.\par
L. Stoyanov, Spectra of Ruelle transfer operators for Axiom A flows, arXiv:0810.1126v4, 22 October 2010. Checked in particular Theorem 1.1, Corollaries 1.4–1.5, and LNIC.\par
T. Humbert, First Ruelle resonance for an Anosov flow with smooth potential, arXiv:2402.04948v1, 7 February 2024. The title is that of the actually acquired v1 PDF.\par
S. Dyatlov and C. Guillarmou, Pollicott–Ruelle resonances for open systems, arXiv:1410.5516v2, 3 April 2016. Checked the hypotheses, resolvent, projector and wavefront statements, and Proposition 6.2.\par
L. Jin and Z. Tao, Counting Pollicott–Ruelle resonances for Axiom A flows, arXiv:2306.02297v2, 8 January 2024.\par
S. Dyatlov and C. Guillarmou, Afterword: Dynamical zeta functions for Axiom A flows, arXiv:1801.00348v2, 31 January 2018.\par
A. Sambarino, Quantitative properties of convex representations, arXiv:1104.4705v2, 28 January 2012. Checked Theorem 3.2 and its reparametrization proof in the stated scope.\par
O. Guichard and A. Wienhard, Anosov representations: Domains of discontinuity and applications, arXiv:1108.0733v2, 31 January 2012. Checked Definition 2.10, Remark 2.11, and Proposition 4.3.\par
O. Glorieux, D. Monclair, and N. Tholozan, Hausdorff dimension of limit sets for projective Anosov representations, J. Éc. polytech. Math. 10 (2023), 1157–1193, DOI: 10.5802/jep.241. The acquired PDF has a cover page; for its text pages, printed page equals physical page plus 1155.\par
J. Danciger, F. Guéritaud, and F. Kassel, Convex cocompact actions in real projective geometry, arXiv:1704.08711v4, 18 April 2023. Checked Definition 1.23, Theorem 1.24, and Proposition 8.1 with its proof.\par
Y. Borns-Weil and S. Shen, Dynamical zeta functions in the nonorientable case, arXiv:2007.08043v2, 1 September 2021. Checked the local basic-set statement and orientation-twisted trace proof.\par
M. Bridgeman, R. Canary, F. Labourie, and A. Sambarino, The pressure metric for Anosov representations, arXiv:1301.7459v4, 1 February 2015. Checked Theorem 1.10, the p. 9 non-elementary convention, and Propositions 4.1–4.2 with their proof. M. Pollicott, On the rate of mixing of Axiom A flows, Invent. Math. 81 (1985), 413–426, DOI: 10.1007/BF01388579. The article-specific Göttingen scan has one cover page; printed text page equals physical page plus 411. Theorem 2 and its proof were checked in the original scan.\par
D. Dolgopyat and M. Pollicott, Addendum to Periodic orbits and dynamical spectra, Ergodic Theory Dynam. Systems 18 (1998), 293–301, DOI: 10.1017/S0143385798113913. Actually acquired: the ten-page author version hosted by Dolgopyat’s University of Maryland mathematics page; its pagination is not the journal pagination.\par
\end{document}
